\def\year{2022}\relax
%File: formatting-instructions-latex-2022.tex
%release 2022.1
\documentclass[letterpaper]{article} % DO NOT CHANGE THIS
\usepackage{aaai22}  % DO NOT CHANGE THIS
\usepackage{times}  % DO NOT CHANGE THIS
\usepackage{helvet}  % DO NOT CHANGE THIS
\usepackage{courier}  % DO NOT CHANGE THIS
\usepackage[hyphens]{url}  % DO NOT CHANGE THIS
\usepackage{graphicx} % DO NOT CHANGE THIS
\urlstyle{rm} % DO NOT CHANGE THIS
\def\UrlFont{\rm}  % DO NOT CHANGE THIS
\usepackage{natbib}  % DO NOT CHANGE THIS AND DO NOT ADD ANY OPTIONS TO IT
\usepackage{caption} % DO NOT CHANGE THIS AND DO NOT ADD ANY OPTIONS TO IT
\DeclareCaptionStyle{ruled}{labelfont=normalfont,labelsep=colon,strut=off} % DO NOT CHANGE THIS
\frenchspacing  % DO NOT CHANGE THIS
\setlength{\pdfpagewidth}{8.5in}  % DO NOT CHANGE THIS
\setlength{\pdfpageheight}{11in}  % DO NOT CHANGE THIS
%
% These are recommended to typeset algorithms but not required. See the subsubsection on algorithms. Remove them if you don't have algorithms in your paper.
\usepackage{algorithm}
\usepackage{algorithmic}
%%%%%%%%%%
% TODO: delete thise
% \usepackage[breaklinks=true,
% citecolor=blue,    
% colorlinks=true,
% linkcolor=red,
% filecolor=magenta,      
% urlcolor=cyan,]{hyperref}
%%%%%%%%%%%

%
% These are are recommended to typeset listings but not required. See the subsubsection on listing. Remove this block if you don't have listings in your paper.
\usepackage{newfloat}
\usepackage{listings}
\lstset{%
	basicstyle={\footnotesize\ttfamily},% footnotesize acceptable for monospace
	numbers=left,numberstyle=\footnotesize,xleftmargin=2em,% show line numbers, remove this entire line if you don't want the numbers.
	aboveskip=0pt,belowskip=0pt,%
	showstringspaces=false,tabsize=2,breaklines=true}
\floatstyle{ruled}
\newfloat{listing}{tb}{lst}{}
\floatname{listing}{Listing}

% add my package 
\usepackage{amsmath}
\usepackage{amsfonts}
\usepackage{pifont}
%
%\nocopyright
%
% PDF Info Is REQUIRED.
% For /Title, write your title in Mixed Case.
% Don't use accents or commands. Retain the parentheses.
% For /Author, add all authors within the parentheses,
% separated by commas. No accents, special characters
% or commands are allowed.
% Keep the /TemplateVersion tag as is
\pdfinfo{
	/Title (Resistance Training using Prior Bias: toward Unbiased Scene Graph Generation)
	/Author (anonymous)
	/TemplateVersion (2022.1)
}

% DISALLOWED PACKAGES
% \usepackage{authblk} -- This package is specifically forbidden
% \usepackage{balance} -- This package is specifically forbidden
% \usepackage{color (if used in text)
% \usepackage{CJK} -- This package is specifically forbidden
% \usepackage{float} -- This package is specifically forbidden
% \usepackage{flushend} -- This package is specifically forbidden
% \usepackage{fontenc} -- This package is specifically forbidden
% \usepackage{fullpage} -- This package is specifically forbidden
% \usepackage{geometry} -- This package is specifically forbidden
% \usepackage{grffile} -- This package is specifically forbidden
% \usepackage{hyperref} -- This package is specifically forbidden
% \usepackage{navigator} -- This package is specifically forbidden
% (or any other package that embeds links such as navigator or hyperref)
% \indentfirst} -- This package is specifically forbidden
% \layout} -- This package is specifically forbidden
% \multicol} -- This package is specifically forbidden
% \nameref} -- This package is specifically forbidden
% \usepackage{savetrees} -- This package is specifically forbidden
% \usepackage{setspace} -- This package is specifically forbidden
% \usepackage{stfloats} -- This package is specifically forbidden
% \usepackage{tabu} -- This package is specifically forbidden
% \usepackage{titlesec} -- This package is specifically forbidden
% \usepackage{tocbibind} -- This package is specifically forbidden
% \usepackage{ulem} -- This package is specifically forbidden
% \usepackage{wrapfig} -- This package is specifically forbidden
% DISALLOWED COMMANDS
% \nocopyright -- Your paper will not be published if you use this command
% \addtolength -- This command may not be used
% \balance -- This command may not be used
% \baselinestretch -- Your paper will not be published if you use this command
% \clearpage -- No page breaks of any kind may be used for the final version of your paper
% \columnsep -- This command may not be used
% \newpage -- No page breaks of any kind may be used for the final version of your paper
% \pagebreak -- No page breaks of any kind may be used for the final version of your paperr
% \pagestyle -- This command may not be used
% \tiny -- This is not an acceptable font size.
% \vspace{- -- No negative value may be used in proximity of a caption, figure, table, section, subsection, subsubsection, or reference
% \vskip{- -- No negative value may be used to alter spacing above or below a caption, figure, table, section, subsection, subsubsection, or reference

\setcounter{secnumdepth}{2} %May be changed to 1 or 2 if section numbers are desired.

% The file aaai22.sty is the style file for AAAI Press
% proceedings, working notes, and technical reports.
%

% Title

% Your title must be in mixed case, not sentence case.
% That means all verbs (including short verbs like be, is, using, and go),
% nouns, adverbs, adjectives should be capitalized, including both words in hyphenated terms, while
% articles, conjunctions, and prepositions are lower case unless they
% directly follow a colon or long dash
\title{Resistance Training using Prior Bias: toward Unbiased Scene Graph Generation}
%Resistance bias for unbiased scene graph generation

% \iffalse
\author{Anonymous Author(s)}
\affiliations{Affiliation }
% \fi

\iffalse
\author{
	%Authors
	% All authors must be in the same font size and format.
% 	Written by AAAI Press Staff\textsuperscript{\rm 1}\thanks{With help from the AAAI Publications Committee.}\\
% 	AAAI Style Contributions by Pater Patel Schneider,
% 	Sunil Issar,\\
% 	J. Scott Penberthy,
% 	George Ferguson,
% 	Hans Guesgen,
% 	Francisco Cruz\equalcontrib,
% 	Marc Pujol-Gonzalez\equalcontrib
}
\affiliations{
	%Afiliations
% 	\textsuperscript{\rm 1}Association for the Advancement of Artificial Intelligence\\
	% If you have multiple authors and multiple affiliations
	% use superscripts in text and roman font to identify them.
	% For example,
	
	% Sunil Issar, \textsuperscript{\rm 2}
	% J. Scott Penberthy, \textsuperscript{\rm 3}
	% George Ferguson,\textsuperscript{\rm 4}
	% Hans Guesgen, \textsuperscript{\rm 5}.
	% Note that the comma should be placed BEFORE the superscript for optimum readability
	
% 	2275 East Bayshore Road, Suite 160\\
% 	Palo Alto, California 94303\\
	% email address must be in roman text type, not monospace or sans serif
% 	publications22@aaai.org
	%
	% See more examples next
}
\fi
%Example, Single Author, ->> remove \iffalse,\fi and place them surrounding AAAI title to use it
\iffalse
\title{My Publication Title --- Single Author}
\author {
	Author Name
}
\affiliations{
	Affiliation\\
	Affiliation Line 2\\
	name@example.com
}
\fi

\iffalse
%Example, Multiple Authors, ->> remove \iffalse,\fi and place them surrounding AAAI title to use it
\title{My Publication Title --- Multiple Authors}
\author {
	% Authors
	First Author Name,\textsuperscript{\rm 1}
	Second Author Name, \textsuperscript{\rm 2}
	Third Author Name \textsuperscript{\rm 1}
}
\affiliations {
	% Affiliations
	\textsuperscript{\rm 1} Affiliation 1\\
	\textsuperscript{\rm 2} Affiliation 2\\
	firstAuthor@affiliation1.com, secondAuthor@affilation2.com, thirdAuthor@affiliation1.com
}
\fi


\begin{document}

% \maketitle

%%%%%%%%%%% start of main context %%%%%%%%%%%%%%%%%%%%%%%%%%
\section{Supplementary Material}
\subsection{Visualization of the loss}
The softmax CE with our resistance bias is calculated as follows:
\begin{align}
    \mathcal{\hat{L}}_i 
	& = -log(\hat{p}_i)
		  = -log \frac{e^{z_{i} - b_{i}}}{\sum_{j\in \mathcal{C}_r}e^{z_{j} - b_{j}}}
    \nonumber\\
    & = -log \frac{e^{z_i}}{\sum_{k\in \mathcal{C}_r}e^{z_k}} \frac{e^{-b_{i}}{\sum_{k\in \mathcal{C}_r}e^{z_k}}}{\sum_{j\in \mathcal{C}_r}e^{z_j-b_{j}}}\nonumber\\
	& = -log \frac{e^{z_i}}{\sum_{k\in \mathcal{C}_r}e^{z_k}} + b_i + log \sum_{j\in\mathcal{C}_r}\frac{e^{-b_j}e^{z_j}}{\sum_{k\in \mathcal{C}_r}e^{z_k}} 
	\nonumber\\
% 	& = -log \frac{e^{z_i}}{\sum_{j\in \mathcal{C}_r}e^{z_j}} + b_i + log \sum_{j\in\mathcal{C}_r}{e^{-b_j}p_j} 	\nonumber\\
	& = -log~p_i + b_i + log \sum_{j\in\mathcal{C}_r}{e^{-b_j}p_j} \nonumber\\
	& = \mathcal{L}_i + \theta_{i}, \label{eq:biasloss}\nonumber\\
% 	& \theta_{i}  =  b_{i}  + log \sum_{j\in \mathcal{C}_r}e^{-b_{j}}p_j,
\end{align}
where $\mathcal{L}_i$ is the vanilla CE, and $~ \theta_{i}  =  b_{i}  + log \sum_{j\in \mathcal{C}_r}e^{-b_{j}}p_j$ is the extra weight of loss value, which is correlated to the resistance bias value for each type. Specifically, the extra weight $\theta_i$ increase the loss value when the resistance bias is large. 
%这一句话，不知道在说什么
% 以二分类的情况为例，我们可视化了有具体bias值的组合对应的损失函数值。
% 如果你读起来有点读不懂，你帮忙改一下看应该怎么描述。(我觉得没太大问题)
% (因为loss与预测分布相关，所以我们只能可视化二分类，因为p2=1-p1.)

Take binary classification as an example, we can visualize the loss value for the two classes with specific resistance bias values. The loss values for several pair of resistance bias values are shown in Figure \ref{fig:loss}. 
% The resistance bias are set by a specific.
\begin{figure}[ht!]
	\centering
	\includegraphics[width=0.9\columnwidth]{figure/lossdist} % Reduce the figure size so that it is slightly narrower than the column.
	\caption{The loss value for cross-entropy with our RTPB. This figure shows three pairs of loss with different resistance bias. In the pair \ding{172}, the resistance bias is set to be $-log(0.9)$ and $-log(0.1)$. $-log(0.7)$ and $-log(0.1)$ is the bias value for the second pair, \ding{173}. In the last pair, \ding{174}, two classes take the same resistance bias value, and the loss value is equivalent to the vanilla CE \cite{Mannor2005Cross}. }
	\label{fig:loss}
\end{figure}

\subsection{Experiment Results}
\begin{figure}[!ht]
	\centering
	\includegraphics[width=0.95\linewidth]{figure/R_mR_pred.png} % Reduce the figure size so that it is slightly narrower than the column.
	\caption{The distribution of R@100 and mR@100 in PredCls task.}
	\label{fig:R_mR_predcls}
\end{figure}
\begin{figure}[!ht]
	\centering
	\includegraphics[width=0.95\linewidth]{figure/R_mR_sgcls.png} % Reduce the figure size so that it is slightly narrower than the column.
	\caption{The distribution of R@100 and mR@100 in SGCls task.}
	\label{fig:R_mR_sgcls}
\end{figure}
\begin{figure}[!ht]
	\centering
	\includegraphics[width=0.95\linewidth]{figure/R_mR_sgdet.png} % Reduce the figure size so that it is slightly narrower than the column.
	\caption{The distribution of R@100 and mR@100 in SGDet task.}
	\label{fig:R_mR_sgdet}
\end{figure}
Table \ref{tab:sgg-r-mr} shows the Recall results and mean Recall results on the VG dataset \cite{krishna2016visual}.
\begin{table*}[!ht]
	\centering
	% 	\scriptsize
	\small
	\begin{tabular}{c|cc|cc|cc}
		\hline
		& \multicolumn{2}{c|}{Predcls} & \multicolumn{2}{c|}{SGcls} & \multicolumn{2}{c}{SGdet} \\
		Models                         & mR@50/100     & R@50/100     & mR@50/100    & R@50/100    & mR@50/100    & R@50/100    \\ \hline
		KERN\ddag\cite{chen2019knowledge}                      & 17.7/19.2     & 67.6/65.8    & 9.4/10.0     & 36.7/37.4   & 6.4/7.3      & 29.8/27.1   \\
		PCPL\ddag\cite{yan2020pcpl}                      & 35.2/37.8     & 50.8/52.6    & 18.6/19.6    & 27.6/28.4   & 9.5/11.7     & 14.6/18.6   \\
		GPS-Net\ddag\cite{lin2020gps}                        & 21.3/22.8   & 66.9/68.8    & 11.8/12.6      & 39.2/40.1   &8.7/9.8      & 28.4/31.7   \\
		% TDE\dag\cite{Tang2020_unbiased}                        & 23.8/27.7     & 43.9/49.8    & 14.1/16.0    & 24.6/27.0   & 10.9/13.3    & 16.2/19.1   \\
		BGNN\cite{li2021bipartite}                           & 30.4/32.9     & 59.2/61.3    & 14.3/16.5    & 37.4/38.5   & 10.7/12.6    & 31.0/35.8   \\ \hline
		MOTIFS\dag\cite{Zellers_2018_CVPR}                    & 17.2/18.6     & 65.4/67.2    & 9.7/10.3     & 40.3/41.1   & 7.7/9.2      & 33.0/37.7   \\
		MOTIFS+TDE\cite{Tang2020_unbiased}                        & 24.9/28.3     & 50.8/55.8    & 13.9/15.2    & 27.2/29.5   & 8.5/9.9    & 7.4/8.4  \\
		MOTIFS+DLFE\cite{chiou2021recovering}                    & 26.9/28.8     & 52.5/54.2    & 15.2/15.9    & 32.3/33.1   & 11.7/13.8    & 25.4/29.4   \\
		MOTIFS + \textbf{RTPB(CB)}               & 35.3/37.7     & 40.4/42.5    & 20.0/21.0    & 26.0/26.9   & 13.1/15.5    & 19.0/22.5   \\ \hline
		VCTree\dag\cite{VCTree_tang}                    & 16.8/18.1     & 65.8/67.5    & 10.8/11.5    & 45.4/46.3   & 7.4/8.6      & 31.0/35.0   \\
		VCTree+TDE\cite{Tang2020_unbiased}                        & 23.3/26.6     & 49.9/54.5    & 11.8/13.4 & 28.8/31.2   & 8.6/10.3    & 19.6/23.3  \\
		VCTree + DLFE\cite{chiou2021recovering}                  & 25.3/27.1     & 51.8/53.5    & 18.9/20.0    & 33.5/34.6   & 11.8/13.8    & 22.7/26.3   \\
		VCTree+\textbf{RTPB(CB)  }               & 33.4/35.6     & 41.2/43.3    &\textbf{ 24.5/25.8}    & 28.7/30.0   & 12.8/15.1    & 18.1/21.3   \\ \hline
		DTrans                         & 19.3/21.0     & 64.4/66.4    & 12.1/13.0    & 39.4/40.3   & 9.0/10.8     & 32.1/36.5   \\
		DTrans+\textbf{RTPB(CB) }                & \textbf{36.2/38.1}    & 45.6/47.5    & 21.8/22.8    & 24.5/25.5   & \textbf{16.5/19.0}    & 19.7/23.4   \\
		DTrans+\textbf{RTPB(VB)}   & 31.1/33.3                    & 56.5/58.4                  & 20.1/21.3                 & 33.6/34.5 & 15.7/18.4 & 27.3/31.4 \\
		DTrans+\textbf{RTPB(PB)}         & 21.6/23.1        & 56.2/58.0       & 14.0/14.7    & 32.6/33.4   & 10.1/11.9    & 24.6/28.9   \\
		DTrans+\textbf{RTPB(EB)}                 & 26.7/28.4     & 50.6/52.6    & 17.4/18.2    & 27.0/28.0   & 14.1/16.1    & 19.1/22.9   \\
		
		% DTrans+RTPB(M)                 & 34.7/36.9     & 51.9/53.8    & 20.8/22.0    & 30.0/30.9   & 16.2/18.8    & 24.5/28.4   \\ \hline
		\hline
	\end{tabular}
	%	\small
	\caption{ The performance on VG~\cite{krishna2016visual} under graph constraints setting. \dag~indicates the results reproduced using the code of \cite{Tang2020_unbiased}. \ddag~models are with VGG16 backbone \cite{simonyan2014very}, while others are with ResNeXt-101-FPN backbone \cite{Lin_2017_ICCV}. (CB), (VB), (PB) and (EB) indicate the count resistance bias, the valid resistance bias, the pair resistance bias, and the estimated resistance bias, respectively.}
	\label{tab:sgg-r-mr}
\end{table*}
To clearly compare all methods, we visualize the results on the PredCls task, SGCls task, and SGDet task in Figure \ref{fig:R_mR_predcls}, \ref{fig:R_mR_sgcls} and \ref{fig:R_mR_sgdet}, respectively.  As shown in Figure \ref{fig:R_mR_predcls}, \ref{fig:R_mR_sgcls} and \ref{fig:R_mR_sgdet}, models trained with RTPB achieve higher mean recall than previous state-of-the-art methods. Besides, the model trained with different $a$ for RTPB provide examples for head-tail trade-off in SGG.

Moreover, in addition to Recall and mean Recall result with graph constraint, we also report the results without graph constraint of our methods in Table \ref{tab:sgg-r-mr-ng}. As we can observe, RTPB also significantly improves the mean Recall of corresponding methods. Besides, DTrans+RTPB(CB) performs the best on nearly all mean Recall. The Recall is dropped by using RTPB. This is because the less frequently seen relationships have been highlighted, where most ground truths are frequently seen. Therefore, it is commonly seen in unbiased SGG strategies that an increase of mean Recall would be followed with a decrease of Recall \cite{Tang2020_unbiased}. 

\begin{table*}[!ht]
	\centering
	% 	\scriptsize
	\small
	\begin{tabular}{c|cc|cc|cc}
\hline
                                        & \multicolumn{2}{c|}{Predcls}   & \multicolumn{2}{c|}{SGcls}     & \multicolumn{2}{c}{SGdet}     \\
Models                                  & mR@50/100          & R@50/100  & mR@50/100          & R@50/100  & mR@50/100          & R@50/100  \\ \hline
MOTIFS\dag\cite{Zellers_2018_CVPR}      & 33.8/46.0          & 81.7/88.9 & 19.9/26.8          & 50.3/54.0 & 13.2/17.9          & 37.2/44.1 \\
MOTIFS+TDE\S~\cite{Tang2020_unbiased}      & 29.0/38.2          & -         & 16.1/21.1          & -         & 11.2/14.9          & -         \\
MOTIFS+DLFE\S~\cite{chiou2021recovering}   & 30.0/45.8          & -         & 25.6/32.0          & -         & 18.1/23.0          & -         \\
MOTIFS + \textbf{RTPB(CB)}              & 48.0/58.9          & 62.8/74.9 & 27.2/33.3          & 39.6/46.5 & 17.5/22.5          & 25.7/32.6 \\ \hline
VCTree\dag\cite{VCTree_tang}            & 34.6/46.9          & 82.5/89.5 & 22.2/30.7          & 56.3/60.6 & 12.6/16.8          & 33.6/39.9 \\
VCTree+TDE\S~\cite{Tang2020_unbiased}      & 32.4/41.5          & -         & 19.1/25.5          & -         & 11.5/15.2          & -         \\
VCTree + DLFE\S~\cite{chiou2021recovering} & 44.6/56.8          & -         & 31.4/38.8          & -         & 17.5/22.5          & -         \\
VCTree+\textbf{RTPB(CB)  }              & 46.4/57.8          & 62.6/74.9 & \textbf{34.0/41.1} & 45.1/53.7 & 17.3/21.9          & 24.5/31.4 \\ \hline
DTrans                                  & 36.6/48.8          & 80.5/87.9 & 23.2/29.7          & 49.5/53.3 & 14.4/19.6          & 36.3/43.0 \\
DTrans+\textbf{RTPB(CB) }               & \textbf{52.6/63.5} & 70.0/80.7 & 31.6/37.1          & 42.0/48.5 & \textbf{22.4/27.5} & 27.5/35.3 \\
DTrans+\textbf{RTPB(VB)}                & 47.5/59.7          & 76.5/85.3 & 30.6/36.9          & 46.6/51.4 & 21.3/27.2          & 32.5/39.8 \\
DTrans+\textbf{RTPB(PB)}                & 38.6/50.4          & 79.2/87.7 & 25.4/31.9          & 48.4/52.8 & 16.6/21.7          & 33.4/41.1 \\
DTrans+\textbf{RTPB(EB)}                & 42.5/52.6          & 76.4/85.9 & 27.5/33.0          & 46.1/51.6 & 19.9/24.7          & 30.4/38.2 \\ \hline
\end{tabular}
	%	\small
	\caption{ The performance on VG~\cite{krishna2016visual} without graph constraint. \dag, \ddag, (CB), (VB), (PB) and (EB) are with the same meaning as in Table \ref{tab:sgg-r-mr}. $\S~$ denote the results are reported by \cite{chiou2021recovering} 
	}
	\label{tab:sgg-r-mr-ng}
\end{table*}
% By 
\subsection{Cost-sensitive Methods}
% In this section, we compare our methods with several typical cost-sensitive methods for imbalance learning. We also demonstrate the performance of different variants of resistance bias.
% \section{Cost-Sensitive Methods}
\begin{table*}[!ht]
    \centering
	\begin{tabular}{c|cc|cc|cc}
		\hline
		& \multicolumn{2}{c|}{Predcls} & \multicolumn{2}{c|}{SGcls} & \multicolumn{2}{c}{SGdet} \\
		Models                         & mR@50/100     & R@50/100     & mR@50/100    & R@50/100    & mR@50/100    & R@50/100    \\ \hline
		Baseline   & 19.3/21.0     & 64.4/66.4    & 12.1/13.0    & 39.4/40.3   & 9.0/10.8     & 32.1/36.5           \\
		$\mathcal{L}_{Reweight}$  & 35.1/\textbf{38.8} & 26.8/30.6 & 14.3/16.8  & 13.6/15.6 & 2.0/2.6 & 4.0/5.1           \\
		$\mathcal{L}_{ClsBal}$     & 19.9/21.9 & 64.5/66.3 & 15.5/17.0  & 38.8/39.7 & 10.7/12.8 & 31.8/36.2           \\
		$\mathcal{L}_{focal}$     & 19.9/21.7 & 63.2/65.6 & 11.3/12.1  & 38.6/39.8 & 8.3/10.0 & 29.6/33.9          \\
		$\mathcal{L}_{LDAM}$      & 15.3/16.6 & 65.1/66.9 & 10.3/10.8    & 39.6/40.4  & 7.0/8.6 & 32.3/37.1        \\
		RTPB(CB)      & \textbf{36.2}/38.1    & 45.6/47.5    & \textbf{21.8/22.8}    & 24.5/25.5   & \textbf{16.5/19.0}    & 19.7/23.4 \\ \hline
	\end{tabular}
	\caption{Comparison of Recall rate and mean Recall rate with other cost-sensitive methods.}
	\label{tab:debias}
\end{table*}
\begin{table*}[!ht]
    \centering
\begin{tabular}{c|cc|cc|cc}
\hline
                         & \multicolumn{2}{c|}{Predcls} & \multicolumn{2}{c|}{SGcls} & \multicolumn{2}{c}{SGdet} \\
Models                   & mR@50/100     & R@50/100     & mR@50/100    & R@50/100    & mR@50/100    & R@50/100    \\ \hline
Baseline                 & 19.3/21.0     & 64.4/66.4    & 12.1/13.0    & 39.4/40.3   & 9.0/10.8     & 32.1/36.5   \\
$\mathcal{L}_{Reweight}$ & 41.2/50.0     & 40.4/53.0    & 16.4/21.1    & 20.3/26.7   & 2.7/3.7      & 5.4/7.0     \\
$\mathcal{L}_{ClsBal}$   & 38.4/51.5     & 80.3/87.9    & 26.9/34.7    & 49.0/52.9   & 16.9/22.3    & 35.9/42.7   \\
$\mathcal{L}_{focal}$    & 27.0/36.9     & 76.6/84.8    & 18.4/24.2    & 46.9/51.4   & 12.4/16.8    & 33.2/39.5   \\
$\mathcal{L}_{LDAM}$     & 33.6/44.9     & 80.2/87.9    & 22.6/28.9    & 49.0/52.9   & 13.0/17.9    & 34.8/42.1   \\
RTPB(CB)                 & \textbf{52.6/63.5}     & 70.0/80.7    & \textbf{31.6/37.1}    & 42.0/48.5   & \textbf{22.4/27.5}    & 27.5/35.3   \\ \hline
\end{tabular}
	\caption{Comparison of Recall rate and mean Recall rate with other cost-sensitive methods.}
	\label{tab:debias_ng}
\end{table*}
We compare our RTPB with the following cost-sensitive methods. 

1) \textbf{Re-weighting Loss}: This strategy uses the fraction of the count of each class to weigh the loss value as follows,
\begin{equation}
	\mathcal{L}_{Reweight}(i) = \frac{1}{n_i}\mathcal{L}_i, 
\end{equation}
where $n_i$ is the number of all samples of $i$-th class in training set. For SGG, $n_i$ is set to be the number of relationship $i$.


2) \textbf{Class Balanced Re-weighting Loss} \cite{cui2019class}: This method uses the inversed effective number of samples to re-balance the loss. The loss is calculated as
\begin{equation}
	\mathcal{L}_{ClsBal}(i) = \frac{1}{n_i^{eff}}\mathcal{L}_i = \frac{1-\beta}{1-\beta^{n_i}}\mathcal{L}_i,
\end{equation}
where $n_i^{eff} = \frac{1-\beta}{1-\beta^{n_i}}$ is the effective number for class $i$, $n_i$ is the number of sample of class $i$. Following the setting of  \cite{cui2019class}, we set the hyper-parameter $\beta= 0.999$. 

% We evaluate different value of $\beta$ on the SGcls task, and the results are shown in 
3) \textbf{Focal Loss} \cite{Lin_2017_ICCV}: This method uses the focal weight to adjust the loss value for well-learned samples and focuses on the hard ones.
\begin{equation}
	\mathcal{L}_{focal}(p_i) = -\alpha(1-p_i)^\gamma log(p_i),
\end{equation}
where $p_i$ is the classification probability for class i. We followed the hyper-parameters ($\gamma = 2.0, \alpha= 0.25$). 

4) \textbf{Label Distribution-Aware Margin Loss} \cite{cao2019learning}: This method uses the margin for each class to balance the model. The loss value for class $i$ is
\begin{equation}
	\begin{split}
		\mathcal{L}_{LDAM}(i) = -log \frac{e^{z_i - \Delta_i}}{e^{z_i - \Delta_i} + \sum_{j\neq i}{e^{z_j}}} \\
		where ~ \Delta_i = \frac{C}{n_i^{1/4}}
	\end{split}
\end{equation}
where $C$ is the constant for the margin size and we set it as 0.5 same as \cite{cao2019learning}, and  $n_i$ is the number of sample of the corresponding class $i$ in the training set. This class-related loss can be regarded as a special form of our method if we set the bias corresponding to the ground truth as follow
\begin{equation}
	b_i =\left\{
	\begin{aligned}
		\Delta_i, ~ & i = gt   \\
		0, ~        & i\neq gt
	\end{aligned}
	\right.
	~,
\end{equation}
%(i.e. re-weighting, focal loss, and margin loss). 
%The re-weighting method.
where $gt$ is the groundtruth of relationship.
We evaluate methods mentioned above on the VG data set\cite{krishna2016visual} with our DTrans, and find out that our RTPB outperforms the rest of the methods by a clear margin in the mean recall metric. Full results are shown in the Table \ref{tab:debias}.

\subsection{Qualitative Study}
We visualized several SGCls examples generated by the DTrans and the DTrans with RTPB(CB) in Figure \ref{fig:sample}. As shown in the figure, the model without RTPB prefers generalized descriptions like $on$ and $has$, which provide less valuable information. This is because these generalized relationships or the frequently seen relationships dominate the training set. However, the model trained with RTPB tends to replace these rough-grained relationships with fine-grained and more informative ones like $parked~on$, $mounted~on$, and $painted~on$, which constitutes more meaningful scene graphs. The differences between the results from the two models demonstrate the necessity of unbiased scene graph generation and the effectiveness of our proposed RTPB.
\begin{figure*}[!ht]
	\centering
	\includegraphics[width=\linewidth]{figure/sample.pdf} % Reduce the figure size so that it is slightly narrower than the column.
	\caption{This figure shows several sample of SGCls. For each sample, we present the groundtruth of object detection in the image and two scene graphs generated by different methods. The scene graphs in the blue boxes are produced by our DTrans without RTPB, and the ones in the red boxes are produced by the DTrans trained with RTPB(CB). We highlight the differences between each pair of scene graphs. Because of the space limitation, we exclude part of trivial detected objects in generated scene graphs. }
	\label{fig:sample}
\end{figure*}
% end of paper content %%
\bibliography{ref}
\end{document}
